% Generated draft: requires mathematical and source review.
\documentclass[11pt]{article}
\usepackage[margin=25mm]{geometry}
\usepackage{fontspec}
\setmainfont{Latin Modern Roman}
\usepackage{amsmath,amssymb,mathtools,unicode-math}
\setmathfont{Latin Modern Math}
\usepackage{enumitem,needspace,xurl,hyperref}
\setlength{\parindent}{0pt}
\setlength{\parskip}{5pt}
\setlength{\emergencystretch}{3em}
\newcommand{\N}{\mathbb{N}}
\newcommand{\Z}{\mathbb{Z}}
\newcommand{\Q}{\mathbb{Q}}
\newcommand{\R}{\mathbb{R}}
\newcommand{\C}{\mathbb{C}}
\newcommand{\F}{\mathbb{F}}
\newcommand{\A}{\mathbb{A}}
\newcommand{\PP}{\mathbb{P}}
\newcommand{\E}{\mathbb{E}}
\newcommand{\eps}{\varepsilon}
\newcommand{\dd}{\,\mathrm{d}}
\begin{document}
\section*{UnsolvedMath ID: 700318}
\textbf{Source problem code:} catalogue entry
\subsection{Definitions and mathematical statement}
Let \(F\) be a field and let \(V\) be an \(n\)-dimensional vector space over \(F\). Write \(\operatorname{GL}(V)\) for the group of invertible \(F\)-linear maps \(V\to V\), and write \(I_V\) for the identity map. A transvection is an element \(t\in\operatorname{GL}(V)\) such that
\[
\operatorname{rank}(t-I_V)=1,\qquad (t-I_V)^2=0.
\]
Its residual line is \(\operatorname{im}(t-I_V)\), and its fixed space is \(\ker(t-I_V)\). A subgroup \(G\leq\operatorname{GL}(V)\) is rich in transvections if \(n\geq2\) and, for every hyperplane \(H\leq V\) and every one-dimensional subspace \(L\leq H\), there exists a transvection \(t\in G\) with
\[
\operatorname{im}(t-I_V)=L,\qquad \ker(t-I_V)=H.
\]
The source asks to describe the automorphisms of subgroups of \(\operatorname{GL}_2(V)\) rich in transvections. Under the intended two-dimensional reading, the problem is: for every field \(F\), every two-dimensional \(F\)-vector space \(V\), and every subgroup \(G\leq\operatorname{GL}(V)\) rich in transvections, describe all abstract group automorphisms of \(G\), namely all bijections \(\alpha:G\to G\) satisfying
\[
\alpha(gh)=\alpha(g)\alpha(h)\qquad(g,h\in G).
\]
In dimension two, the richness condition is equivalently that every line \(\ell\leq V\) occurs as both the residual line and the fixed space of some transvection in \(G\). No restriction on the characteristic or cardinality of \(F\) is supplied. The notation \(\operatorname{GL}_2(V)\) is retained here as a source qualification: interpreting it as the linear group on a two-dimensional space, and interpreting automorphisms as abstract group automorphisms, require confirmation in mathematical review.
\subsection{Short English statement}
For each field and each two-dimensional vector space, describe all group automorphisms of every linear subgroup containing a transvection with each prescribed line as its fixed space and residual line. This uses the two-dimensional interpretation of the source notation; the supplied registry labels the problem open.
\subsection{Sources}
\begin{itemize}
\item Local UnsolvedMath registry, permanent ID \(700318\), exact source problem code \texttt{catalogue entry}, title \emph{Kourovka Notebook Problem 5.35}. The supplied record labels the problem open.
\item \emph{The Kourovka Notebook (Unsolved Problems in Group Theory)}, fifth issue (1976), Problem 5.35; proposed by Yu. I. Merzlyakov.
\item The supplied background also identifies the twenty-first issue edition (2026), arXiv:1401.0300v48. The supplied collection description names E. I. Khukhro and V. D. Mazurov as editors. These bibliographic details are supplied attribution and have not been independently verified here.
\end{itemize}
\end{document}
